\hypertarget{ejercicio-3.-documentaciuxf3n-de-los-algoritmos}{%
\section{Ejercicio 3. Documentación de los algoritmos}\label{ejercicio-3.-documentaciuxf3n-de-los-algoritmos}}

En las estimaciones, n = \textbar Q\textbar, k = \textbar Σ\textbar{} y m es el número de transiciones (m ≤ n²(k + 1) en un AFN-λ).

\hypertarget{cerradura-ux3bb-y-simulaciuxf3n-directa}{%
\subsection{Cerradura λ y simulación directa}\label{cerradura-ux3bb-y-simulaciuxf3n-directa}}

\textbf{Problema.} Dado un AFN-λ y un conjunto de estados S, hallar los estados alcanzables desde S sin consumir símbolos; con ello, decidir si el autómata acepta una palabra sin convertirlo.

\textbf{Resultado teórico.} δ̂(q, λ) = λ-cl(q) y δ̂(q, ωa) = λ-cl(⋃\_\{p ∈ δ̂(q, ω)\} δ(p, a)); ω se acepta si δ̂(q₀, ω) ∩ F ≠ ∅ \autocite[sec.~2.5.3--2.5.4]{hopcroft2007}.

\begin{Shaded}
\begin{Highlighting}[]
\NormalTok{CIERRE(S)                                  \# Automata.cierre}
\NormalTok{    visitados ← S; pila ← S}
\NormalTok{    mientras pila ≠ ∅: q ← pop(pila)}
\NormalTok{        para r ∈ δ(q, λ) − visitados: visitados ← visitados ∪ \{r\}; push(pila, r)}
\NormalTok{    devolver visitados}

\NormalTok{SIMULAR(ω)                                 \# Simulacion.avanzar / completar}
\NormalTok{    S ← CIERRE(\{q₀\})}
\NormalTok{    para cada símbolo a de ω: S ← CIERRE(⋃\_\{q∈S\} δ(q, a))     \# Automata.transitar}
\NormalTok{    aceptar si S ∩ F ≠ ∅                                       \# Automata.es\_aceptacion}
\end{Highlighting}
\end{Shaded}

\textbf{Ejemplo (196).} A mano: λ-cl(A) = \{A, B\} porque A →λ B; λ-cl(B) = \{B\}, λ-cl(C) = \{C\}. Con ω = \texttt{01}: inicio \{A, B\}; con 0, δ(A, 0) ∪ δ(B, 0) = \{A\} ∪ \{A, B\} = \{A, B\}, cuyo cierre es \{A, B\}; con 1, δ(B, 1) = \{C\}; C ∈ F, se acepta. El software da la misma tabla de cerraduras y la misma traza.

\textbf{Pasos.} Una cerradura revisa cada estado y cada arista λ una vez: O(n + m\_λ). Cada paso de la simulación une hasta n conjuntos de hasta n estados y los cierra: O(n²), y una palabra cuesta O(\textbar ω\textbar·n²).

\hypertarget{eliminaciuxf3n-de-transiciones-ux3bb}{%
\subsection{Eliminación de transiciones λ}\label{eliminaciuxf3n-de-transiciones-ux3bb}}

\textbf{Problema.} Dado un AFN-λ, obtener un AFN sin transiciones λ con el mismo lenguaje y los mismos estados.

\textbf{Resultado teórico.} Con Δ'(q, a) = λ-cl(δ(λ-cl(q), a)) y F' = \{q : λ-cl(q) ∩ F ≠ ∅\}, el AFN (Q, Σ, Δ', q₀, F') reconoce el mismo lenguaje \autocites[sec.~2.5.5]{hopcroft2007}[sec.~1.2, demostración del teorema 1.39]{sipser2013}.

\begin{Shaded}
\begin{Highlighting}[]
\NormalTok{QUITAR{-}λ(N)                                \# QuitarLambda.resolver}
\NormalTok{    para q ∈ Q: C[q] ← CIERRE(\{q\})          \# etapa 1}
\NormalTok{    para q ∈ Q, a ∈ Σ: Δ\textquotesingle{}(q, a) ← CIERRE(⋃\_\{p∈C[q]\} δ(p, a))   \# etapa 2}
\NormalTok{    F\textquotesingle{} ← \{q : C[q] ∩ F ≠ ∅\}                 \# etapa 3}
\NormalTok{    devolver (Q, Σ, Δ\textquotesingle{}, q₀, F\textquotesingle{})             \# etapa 4}
\end{Highlighting}
\end{Shaded}

\textbf{Ejemplo (196).} A mano: Δ'(A, 0) = λ-cl(δ(\{A, B\}, 0)) = λ-cl(\{A, B\}) = \{A, B\}, que añade B a δ(A, 0) = \{A\}; Δ'(A, 1) = λ-cl(δ(\{A, B\}, 1)) = \{C\}, que no existía. B y C conservan sus transiciones, y como ninguna cerradura distinta de la de C contiene a C, F' = \{C\}. El software marca exactamente esos dos destinos añadidos («(+B)» y «(+C)») y reporta que ningún estado nuevo pasa a ser final.

\textbf{Pasos.} n cerraduras más, para cada uno de los n·k pares, una unión de hasta n conjuntos y un cierre: O(k·n³) en el peor caso.

\hypertarget{construcciuxf3n-de-subconjuntos}{%
\subsection{Construcción de subconjuntos}\label{construcciuxf3n-de-subconjuntos}}

\textbf{Problema.} Dado un AFN (o AFN-λ), obtener un AFD equivalente construyendo sólo los subconjuntos alcanzables.

\textbf{Resultado teórico.} Para todo AFN existe un AFD equivalente cuyos estados son subconjuntos de Q \autocites[teorema 11]{rabin1959}[sec.~2.3.5, teorema 2.11]{hopcroft2007}.

\begin{Shaded}
\begin{Highlighting}[]
\NormalTok{SUBCONJUNTOS(N)                            \# Subconjuntos.resolver}
\NormalTok{    S₀ ← CIERRE(\{q₀\}); descubiertos ← [S₀]; cola ← [S₀]}
\NormalTok{    mientras cola ≠ ∅:}
\NormalTok{        S ← desencolar(cola)                \# una etapa por fila}
\NormalTok{        para a ∈ Σ:}
\NormalTok{            T ← CIERRE(⋃\_\{q∈S\} δ(q, a))     \# Automata.transitar}
\NormalTok{            si T ∉ descubiertos: agregar T a descubiertos y a cola}
\NormalTok{            δ\_D(S, a) ← T}
\NormalTok{    F\_D ← \{S ∈ descubiertos : S ∩ F ≠ ∅\}}
\end{Highlighting}
\end{Shaded}

\textbf{Ejemplo (192).} A mano: A = \{A\}; con 0 va a \{B\} = B y con 1 a \{B, C\} = C. B: con 0 a \{A, B\} = D, con 1 a ∅. C: con 0 a δ(B, 0) ∪ δ(C, 0) = \{A, B, C\} = E. D: con 0 a \{A, B\} (D), con 1 a \{B, C\} (C). E: con 0 a E y con 1 a C. Son finales C y E porque contienen al estado final C. El software genera las cinco filas en el mismo orden y con los mismos nombres; JFLAP (\emph{Convert to DFA}) da un AFD de cinco estados equivalente.

\textbf{Pasos.} A lo sumo 2ⁿ subconjuntos, y por cada uno k transiciones de costo O(n²): O(k·n²·2ⁿ) en el peor caso; en los casos de prueba se generaron entre 2 y 12 subconjuntos.

\hypertarget{minimizaciuxf3n}{%
\subsection{Minimización}\label{minimizaciuxf3n}}

\textbf{Problema.} Dado un AFD, obtener el AFD con el menor número de estados que reconoce el mismo lenguaje.

\textbf{Resultado teórico.} El AFD mínimo es único salvo el nombre de los estados y sus estados corresponden a las clases de la relación de Myhill-Nerode \autocite{nerode1958}; dos estados son equivalentes si ninguna palabra los distingue, y el algoritmo de llenado de tabla encuentra exactamente los pares distinguibles \autocite[sec.~4.4.1--4.4.3]{hopcroft2007}.

\begin{Shaded}
\begin{Highlighting}[]
\NormalTok{MINIMIZAR(D)                               \# Minimizacion.resolver}
\NormalTok{    eliminar estados inaccesibles           \# Automata.alcanzables}
\NormalTok{    si δ es parcial: completar con un sumidero      \# AFD.completar}
\NormalTok{    marcar \{p, q\} si exactamente uno está en F       \# ronda 0}
\NormalTok{    repetir (ronda r = 1, 2, …):}
\NormalTok{        para cada par \{p, q\} sin marcar y a ∈ Σ:}
\NormalTok{            si \{δ(p, a), δ(q, a)\} está marcado: marcar \{p, q\} en la ronda r}
\NormalTok{    hasta que una ronda no marque nada}
\NormalTok{    clases ← pares sin marcar                        \# Minimizacion.\_clases}
\NormalTok{    construir el AFD de clases (y retirar la del sumidero si se agregó)}
\end{Highlighting}
\end{Shaded}

\textbf{Ejemplo (156).} A mano: el único final es E, así que la ronda 0 marca los cuatro pares con E. En la ronda 1, con el símbolo 1, A y C van a C mientras que B y D van a E; como \{C, E\} está marcado, se marcan \{A, B\}, \{A, D\}, \{B, C\} y \{C, D\}. La ronda 2 no marca nada. Quedan \{A, C\} y \{B, D\} sin marcar: el mínimo tiene tres estados, \{A, C\}, \{B, D\} y E. El software muestra la tabla triangular de las tres rondas con las mismas marcas y el mismo resultado.

\textbf{Equivalencia.} Dos AFD son equivalentes si sus mínimos coinciden salvo el nombre; \texttt{Equivalencia} lo decide recorriendo ambos mínimos en paralelo, y si no lo son, una búsqueda en anchura sobre el producto da la palabra más corta que los distingue \autocite[sec.~4.4.2]{hopcroft2007}.

\textbf{Pasos.} Hay n(n − 1)/2 pares y a lo sumo n − 2 rondas útiles, cada una de costo O(k·n²): O(k·n³) en total. El algoritmo de Hopcroft baja el costo a O(k·n log n) \autocite{hopcroft1971}. La prueba de isomorfismo cuesta O(k·n) y la del contraejemplo O(k·n₁·n₂).

\hypertarget{validaciuxf3n-por-lotes}{%
\subsection{Validación por lotes}\label{validaciuxf3n-por-lotes}}

\textbf{Problema.} Validar muchas palabras con un mismo autómata y conservar el resultado. No requiere teoría nueva: es SIMULAR(ω) para cada línea.

\begin{Shaded}
\begin{Highlighting}[]
\NormalTok{LOTE(N, texto)                              \# Lote.validar\_texto}
\NormalTok{    para cada línea no vacía que no empiece con \#:}
\NormalTok{        ω ← (λ si la línea es «λ»; si no, la línea)}
\NormalTok{        registrar (número, ω, ACEPTADA/RECHAZADA/ERROR, estados finales)}
\NormalTok{    escribir informe                         \# Lote.informe → Informe.guardar}
\end{Highlighting}
\end{Shaded}

\textbf{Ejemplo.} Con el AFN 263 (palabras que terminan en 01) y \texttt{cadenas-ejemplo.txt}: \texttt{01}, \texttt{1101}, \texttt{001} y \texttt{0101010101} se aceptan; \texttt{0110}, λ y \texttt{10} se rechazan, y \texttt{012} produce un error porque 2 ∉ Σ. Es lo que se obtiene a mano al revisar el sufijo de cada palabra.

\textbf{Pasos.} La suma, sobre las palabras, de O(\textbar ω\textbar·n²).
